\documentclass[12pt]{article}
\usepackage[T1]{fontenc}
\usepackage[utf8]{inputenc}
\usepackage{lmodern}
\usepackage{textcomp}
\usepackage{enumitem}
\usepackage{geometry}
\geometry{margin=1in}
\begin{document}
\begin{center}
Answer Key - Philosophy - Graduate Level Assessment 14 \
\small (Answers are ordered exactly as in the question paper.)
\end{center}

\section*{Multiple Choice}
\begin{enumerate}[label=\arabic*.]
\item \textbf{Answer:} A
\\ \textit{Options:} A) Correct liberation; B) Correct understanding; C) Correct knowledge; D) Correct faith; E) Correct conduct
\\ \textit{Note:} The term "samyak jnana" in Jaina traditions refers to a comprehensive and correct understanding of reality that leads to liberation from the cycle of birth and death, making "correct liberation" an apt description of its significance.
\item \textbf{Answer:} A
\\ \textit{Options:} A) xPs; B) SPx; C) sP; D) Pxs; E) Spx
\\ \textit{Note:} The correct answer, xPs, accurately captures the assertion that "Sheena" (represented by x) possesses the property of being a "punk rocker" (denoted by the predicate P), thus translating the statement into the appropriate form of predicate logic.
\item \textbf{Answer:} B
\\ \textit{Options:} A) Epistemology; B) Morality; C) Religion; D) Aesthetics
\\ \textit{Note:} In his final work, Laws, Plato emphasizes the importance of moral education and the establishment of a just society, reflecting a shift from cosmological inquiries to a focus on ethical and political issues.
\item \textbf{Answer:} E
\\ \textit{Options:} A) actions that affect only themselves.; B) actions performed under duress.; C) actions that do not violate anyone's rights.; D) actions that benefit others.; E) what is due to factors beyond their control.
\\ \textit{Note:} Nagel argues that individuals should not be held morally accountable for actions influenced by uncontrollable factors, as true moral responsibility requires the ability to choose freely, which is compromised when external circumstances dictate one's behavior.
\item \textbf{Answer:} D
\\ \textit{Options:} A) consequentialism.; B) virtue ethics.; C) utilitarianism.; D) none of the above
\\ \textit{Note:} Kant's moral theory, known as deontological ethics, emphasizes duty and the categorical imperative, distinguishing it from consequentialist theories like utilitarianism and other ethical frameworks, thus making "none of the above" the correct answer.
\end{enumerate}

\section*{True / False}
\begin{enumerate}[label=\arabic*.]
\setcounter{enumi}{5}
\item \textbf{Answer:} True
\\ \textit{Options:} A) True; B) False
\\ \textit{Note:} The statement is true because the straw man fallacy involves misrepresenting an argument to make it easier to attack, and claiming that a person's inability to do good negates their good actions distorts the complexity of moral evaluation, reducing it to a simplistic and inaccurate assertion.
\item \textbf{Answer:} False
\\ \textit{Options:} A) True; B) False
\\ \textit{Note:} Berkeley argues that to exist is to be perceived, meaning that existing and perceiving are intrinsically linked rather than independent, as his philosophy asserts that objects only exist in the minds of perceivers.
\item \textbf{Answer:} False
\\ \textit{Options:} A) True; B) False
\\ \textit{Note:} Financial dependence on partners typically undermines women's autonomy and agency in their sexual choices, rather than empowering them.
\item \textbf{Answer:} False
\\ \textit{Options:} A) True; B) False
\\ \textit{Note:} Singer's principle of equality asserts that the capacity to suffer, rather than species membership, is the basis for moral consideration, thereby rejecting any inherent superiority of humans over nonhuman animals.
\item \textbf{Answer:} True
\\ \textit{Options:} A) True; B) False
\\ \textit{Note:} Marcia Baron argues that the unrealistic assumptions of ticking bomb scenarios undermine their validity, as they simplify complex ethical considerations and distract from the broader implications of justifying torture in real-world situations.
\end{enumerate}

\section*{Long Form Response}
\begin{enumerate}[label=\arabic*.]
\setcounter{enumi}{10}
\item \textbf{Answer:} Valuing a patient's autonomy by maximizing their effective options aligns with Velleman's perspective on informed decision-making, emphasizing the importance of enabling individuals to make choices that reflect their values and preferences. According to Velleman, informed decision-making is not merely about providing information but ensuring that patients can effectively engage with and act upon that information in a way that is meaningful to them. This approach underscores the ethical obligation of healthcare providers to recognize and support the unique circumstances of each patient, thereby fostering an environment where autonomy is respected. By prioritizing the enhancement of patient options, healthcare professionals can facilitate truly informed choices, reinforcing the moral tenet that respects individual agency within the healthcare context. Such practices not only empower patients but also enhance the ethical integrity of healthcare delivery.
\\ \textit{Note:} correct because it captures Velleman's view that informed decision-making in healthcare involves not only providing information but also empowering patients to make choices that genuinely reflect their personal values and preferences, thereby enhancing their autonomy.
\item \textbf{Answer:} The use of a "war" metaphor in discussions of government drug policies significantly obstructs our understanding of the rights of drug users by framing the issue in terms of conflict and criminality rather than individual autonomy and public health. According to Huemer, this militaristic language leads to a dehumanization of drug users, depicting them as enemies rather than individuals with rights and needs. Consequently, it becomes challenging to engage in rational discourse about the ethical implications of drug use and the need for harm reduction strategies. This metaphorical framing inhibits a nuanced understanding of drug users' rights, as it prioritizes punitive measures over compassionate and rights-based approaches to drug policy. As a result, the ongoing "war on drugs" distracts from the potential for reform that respects the dignity and autonomy of individuals who use drugs.
\\ \textit{Note:} correct because it identifies how the "war" metaphor shifts the focus from the rights and needs of drug users to a narrative of conflict, thereby dehumanizing them and complicating ethical discussions about their autonomy and well-being.
\end{enumerate}

\vfill
\noindent\textit{End of Answer Key}
\end{document}